% !TeX root = ../nc-quadric-surface-avatar.tex
% 2026-11-23 / v4.20 / cubic geometric 開の不変式表示。

\subsection{四量子ビット不変式によるgeometric開の二本の補曲線}
\label{subsec:qa-cubic-geometric-open-invariants}

前項の $f=1$ chart と交差縮約の規約を保つ。
固定した Segre 双線型形式を $J$ とし，
\[
 K_T=JT^{\mathrm t}JT,\qquad
 e_i=e_i(K_T)\ (i=1,2,3),\qquad d=\det T
\]
と置く。右辺の $e_i$ は行列の固有値の第 $i$ 基本対称式である。
$e_1,e_2,d,e_3$ の通常の多項式次数は $2,4,4,6$，
前項の指標 $\chi$ による次数は $1,2,2,3$ である。
Okawa--Ueda の不変式環の計算
\cite[\S3, (3.9)--(3.16)]{OU2014} と Newton の恒等式により
\[
 k[T]^{SL_2^4}=k[e_1,e_2,d,e_3].
 \tag{CI}\label{eq:qa-cubic-elementary-invariant-ring}
\]
従って既存 GIT 商には同次座標
$[e_1:e_2:d:e_3]$ を持つ
$\mathbb P(1,2,2,3)$ を用いることができる。

\begin{proposition}[二半不変式の普遍恒等式]
\label{prop:qa-cubic-adjugate-explicit-invariants}
$N(T)=\operatorname{Cr}(\operatorname{adj}T)$，
$p_{\rm adj}=\det N(T)$，
$q_{\rm adj}=e_3(JN(T)^{\mathrm t}JN(T))$ とする。
任意の $4\times4$ 行列 $T$ に対し，次の多項式恒等式が成立する。
\[
 \begin{split}
 16p_{\rm adj}
   &=e_3^2-4d^2e_2+8d^3,\\
 16q_{\rm adj}
   &=e_3^3-4d^2e_2e_3+16d^4e_1-8d^3e_3.
 \end{split}
 \tag{CPQ}\label{eq:qa-cubic-adjugate-pq-formulas}
\]
特に
\[
 q_{\rm adj}-e_3p_{\rm adj}
     =d^3(de_1-e_3).
 \tag{CPQ-diff}\label{eq:qa-cubic-adjugate-pq-difference}
\]
\end{proposition}
\begin{proof}
基底を $e_{00},e_{01},e_{10},e_{11}$ の順とし，
$H$ の列を
\[
 e_{00}+e_{11},\quad e_{00}-e_{11},\quad
 e_{01}+e_{10},\quad e_{01}-e_{10}
\]
とする。
この各列を適当な定数で割れば $J$ に関する正規直交基底となる。
その定数対角行列は以下の対角行列と可換なので，
Cartan 部分空間を
\[
 T=H\operatorname{diag}(x_1,x_2,x_3,x_4)H^{-1}
\]
と書いてよい。この部分空間では
\[
 e_i=e_i(x_1^2,x_2^2,x_3^2,x_4^2),\qquad
 d=x_1x_2x_3x_4 .
\]
$\lambda_i=\prod_{j\ne i}x_j$ と置けば，
$\operatorname{adj}T$ の同じ基底での固有値は $\lambda_i$ である。
交差縮約の係数を直接並べ替えると，$N(T)$ も同じ $H$ で対角化され，
その固有値 $\mu_i$ は
\[
 \begin{pmatrix}\mu_1\\\mu_2\\\mu_3\\\mu_4\end{pmatrix}
 =
 \frac12
 \begin{pmatrix}
 1&1&1&1\\
 1&1&-1&-1\\
 1&-1&1&-1\\
 1&-1&-1&1
 \end{pmatrix}
 \begin{pmatrix}\lambda_1\\\lambda_2\\\lambda_3\\\lambda_4\end{pmatrix}.
 \tag{CH}\label{eq:qa-cross-adjugate-hadamard}
\]
この $N(T)$ は対称で $J$ と可換なので，
$JN(T)^{\mathrm t}JN(T)=N(T)^2$ である。

$E_i=e_i(\lambda_1^2,\ldots,\lambda_4^2)$，
$D=\lambda_1\lambda_2\lambda_3\lambda_4$ と置くと，
式~\eqref{eq:qa-cross-adjugate-hadamard} の四つの線型式の展開により
\[
 \begin{split}
 16\prod_i\mu_i&=E_1^2-4E_2+8D,\\
 16e_3(\mu_1^2,\ldots,\mu_4^2)
     &=E_1^3-4E_1E_2+16E_3-8DE_1.
 \end{split}
\]
一方，補積 $\lambda_i$ の定義から
\[
 E_1=e_3,\qquad E_2=d^2e_2,\qquad
 E_3=d^4e_1,\qquad D=d^3 .
\]
代入すると求める二式を得る。

以上は Cartan 部分空間上の計算だが，両辺は全て $SL_2^4$ 不変式である。
式~\eqref{eq:qa-cubic-elementary-invariant-ring} の直交群による記述では，
Cartan への制限は $W(D_4)$ 不変式環との同型，特に単射である。
従って二式は全行列空間上の普遍恒等式となる。
最後の式は二式を引き算すれば従う。
\end{proof}

\begin{corollary}[Geometric開の被約補集合]
\label{cor:qa-cubic-geometric-two-complement-curves}
$\Omega_{\rm geo}\subset\mathbb P(1,2,2,3)$ を前項の geometric 開の
良モジュライとする。その補集合を被約構造で考えると
\[
 \bigl(\mathbb P(1,2,2,3)\setminus\Omega_{\rm geo}\bigr)_{\rm red}
    =C_0\cup C_1,
\]
\[
 C_0=V(d,e_3),\qquad
 C_1=V(e_3-de_1,\ 4e_2-e_1^2-8d).
 \tag{CC}\label{eq:qa-cubic-two-bad-curves}
\]
両成分は $\mathbb P(1,2)$ と同型であり，
唯一の交点は $[2:1:0:0]$ である。
被約和の同次イデアルは
\[
 \sqrt{(p_{\rm adj},q_{\rm adj})}
 =
 (e_3-de_1,\ d(4e_2-e_1^2-8d)).
 \tag{CR}\label{eq:qa-cubic-geometric-radical}
\]
交点の近傍でこの和は通常の節点である。
\end{corollary}
\begin{proof}
$h=e_3-de_1$，$\ell=4e_2-e_1^2-8d$ と置く。
前命題から，二式のイデアルは単元の係数を除いて
\[
 I=(h^2+2de_1h-d^2\ell,\ d^3h)
\]
と書ける。
これを含む素イデアルは，$d$ を含む場合には $h$ も含み，
$d$ を含まない場合には $h,\ell$ を含む。
逆に $(d,h)$ と $(h,\ell)$ はともに $I$ を含む素イデアルである。
従って
\[
 \sqrt I=(d,h)\cap(h,\ell)=(h,d\ell).
\]
最後の等号は $h$ で割った多項式環において $d$ と $\ell$ が
互いに素であることによる。

$C_0$ の座標環は $k[e_1,e_2]$，$C_1$ の座標環は $k[e_1,d]$ なので，
両者は重み $(1,2)$ の射影直線となる。
例えば後者の閉埋込みは
\[
 [s:t]\longmapsto
 [s:(s^2+8t)/4:t:st]
\]
で与えられる。
共通点では $d=e_3=0$ と $4e_2=e_1^2$ が成り立つため，
射影点は $[2:1:0:0]$ に限られる。

この点は $e_1\ne0$ の滑らかな affine chart に入る。
その座標を
\[
 u=\frac{4e_2}{e_1^2}-1,\qquad
 v=\frac d{e_1^2},\qquad
 w=\frac{e_3}{e_1^3}
\]
とすると，被約和のイデアルは $(w-v,\ v(u-8v))$ である。
滑らかな面 $w=v$ の正則座標 $v,u-8v$ に関する二軸の和なので，
通常の節点である。
\end{proof}

\begin{remark}[被約補集合と検算の範囲]
式~\eqref{eq:qa-cubic-geometric-radical} は radical の等式であり，
二つの高次式が定める閉部分そのものを被約と主張してはいない。
ここでの二曲線は既存 cubic GIT の非 geometric 領域であり，
全 $H_2$ の大域的 good moduli に含まれる noncubic 閉曲線とは区別する。

再現コード \texttt{verify\_cubic\_adjugate\_invariants.py} は，
上の基底と交差縮約を係数から実装し，Cartan 恒等式，
階数四・三・二の非対角行列での直接式，
および radical 証明に使った代数恒等式を exact arithmetic で検算する。
全行列上での恒等式の証明は，これらの有限例ではなく，
上記の不変式制限の単射性によって与えている。
\end{remark}
